\documentclass{article}

\usepackage{neurips_2025}

\usepackage[utf8]{inputenc}
\usepackage[T1]{fontenc}
\usepackage{hyperref}
\usepackage{url}
\usepackage{booktabs}
\usepackage{amsfonts}
\usepackage{amssymb}
\usepackage{amsmath}
\usepackage{nicefrac}
\usepackage{microtype}
\usepackage{xcolor}
\usepackage{graphicx}
\usepackage{algorithm}
\usepackage{algorithmic}
\usepackage{multirow}
\usepackage{subcaption}

\title{Tree-Native Neural Integration: \\
  Graph Message Passing for Symbolic Antidifferentiation}

\author{%
  Anonymous Author(s)
}

\begin{document}

\maketitle

\begin{abstract}
	Neural approaches to symbolic integration flatten expression trees into token
	sequences, discarding the hierarchical structure that defines mathematical
	expressions. We introduce a tree-native architecture for symbolic integration
	that operates directly on expression trees: a graph neural network encoder with
	bidirectional message passing and variable-aware attention, followed by a
	top-down autoregressive tree decoder. The full model has $12.1$M parameters.
	We train on $200$M verified integrand--antiderivative pairs generated by a
	Rust-based backward-construction pipeline, with $40$M pairs per task type spanning five types:
	univariate, multivariate, definite, parametric, and special-function
	integrals. To our knowledge this is the first neural model to produce
	\emph{symbolic} antiderivatives for multivariate integrands,
	$\int f(x,y)\,dx = F(x,y) + g(y)$, where non-integrated variables act as
	symbolic parameters.
	Inference uses grammar-constrained decoding and a sample-and-verify strategy
	that exploits the high-confidence verification oracle provided by
differentiation.
	Trained on $200$M pairs ($40$M per task type), the $12.1$M-parameter
	tree-native model matches the $99.7\%$ univariate solve rate of a
	$95$M-parameter sequence transformer \citep{lample2020deep} and
	achieves $>97\%$ across all five task types, including four not
	addressed by prior neural approaches.
\end{abstract}


%% ====================================================================
\section{Introduction}
\label{sec:intro}
%% ====================================================================

Symbolic integration---finding a closed-form antiderivative $F(x)$ such that
$\tfrac{d}{dx}F(x) = f(x)$---is among the most fundamental tasks in
mathematical reasoning. Computer algebra systems (CAS) implement variants of
the Risch algorithm~\citep{risch1969problem} and heuristic pattern matching,
yet regularly fail on deeply nested compositions, special-function
integrands, and expressions outside their built-in rule
databases~\citep{moses1971symbolic}.

\citet{lample2020deep} demonstrated that encoder-decoder transformers can
learn symbolic integration by treating it as sequence-to-sequence
translation: the integrand is serialized in prefix notation and the
antiderivative is generated token by token. Their model, trained on $40$M
pairs with a $95$M-parameter transformer on $32$ V100 GPUs, reaches
$>\!99\%$ accuracy on held-out univariate integrals. Subsequent work has
refined training data, search strategies, and verification
\citep{davis2019use,alphaintegrator2024,sird2023,meidani2024snip}. All existing neural
approaches, however, share a common representational choice: they flatten
expression trees into linear sequences, discarding the tree structure that
\emph{defines} mathematical expressions.

We argue this flattening is both wasteful and limiting. Mathematical
expressions are trees: operators are internal nodes, operands are leaves,
and the semantics of an expression is determined by its tree structure. A
sequence transformer must \emph{rediscover} this structure from positional
patterns in a flat token stream, spending model capacity on a problem that
is already solved by the input representation.

We propose a \textbf{tree-native} architecture for symbolic integration that
operates directly on expression trees. Our system consists of:
\begin{enumerate}
	\item A \textbf{GNN encoder} with $8$ rounds of bidirectional
	      parent$\leftrightarrow$child message passing, augmented with
	      variable-aware attention that biases nodes sharing the same dependency
	      on the integration variable.
	\item A \textbf{top-down tree decoder} that generates the output
	      expression tree level by level, predicting each node's symbol via
	      cross-attention to the encoder output.
\end{enumerate}
The full model has $12.1$M parameters, roughly an order of magnitude smaller
than the seq2seq models reported by \citet{lample2020deep}. We train on
$200$M integrand--antiderivative pairs generated via backward construction
(differentiate a random antiderivative, pair the derivative as the
integrand) using a Rust pipeline that verifies every pair symbolically. Our
dataset spans five task types with $40$M pairs each: univariate,
multivariate, definite, parametric, and special-function. Prior work on neural multivariate integration produces
\emph{numerical} estimates~\citep{rubab2025neural,maitre2023numerical}; to
our knowledge, no prior neural model predicts \emph{symbolic}
antiderivatives for multivariate integrands
($\int f(x,y)\,dx = F(x,y) + g(y)$). We note that structurally this task
reduces to univariate integration with additional inert parameters; the
key challenge is learning to carry function-valued constants of
integration through the tree decoder.

\paragraph{Contributions.}
\begin{enumerate}
	\item The first tree-native neural architecture for symbolic integration,
	      combining GNN message passing with top-down tree decoding.
	\item The first neural model to produce \emph{symbolic} antiderivatives
	      for multivariate integrands, extending univariate integration with
	      variable-aware attention for inert-parameter handling.
	\item A benchmark spanning five task types with
	      $200$M verified pairs ($40$M per type), exceeding the data scale
	      of prior work while using $8\times$ fewer parameters.
	\item A sample-and-verify inference framework exploiting the
	      high-confidence verification oracle of differentiation, with
	      grammar-constrained decoding to maximize the yield of valid
	      candidates.
\end{enumerate}


%% ====================================================================
\section{Related Work}
\label{sec:related}
%% ====================================================================

\paragraph{Neural symbolic integration.}
\citet{lample2020deep} introduced the sequence-to-sequence approach to
symbolic integration, training transformers on backward-generated data
(differentiate random expressions to obtain integrand--antiderivative
pairs). Their $95$M-parameter model, trained on $40$M pairs with $32$
V100 GPUs, achieves $>\!99\%$ accuracy on a held-out test set of
univariate integrals, establishing the seq2seq approach as the dominant
paradigm~\citep{vaswani2017attention}.
\citet{alphaintegrator2024} augmented this pipeline with
transformer-guided action search for harder integrals.
SIRD~\citep{sird2023} predicts integration \emph{rules} rather than
antiderivatives directly, recovering a form of interpretability.
\citet{davis2019use} critiqued the original seq2seq evaluation, argued
for structural splits, and demonstrated \emph{coefficient brittleness}:
seq2seq models fail on $\int 88x^{42}\,dx$ despite solving
$\int 26x^{42}\,dx$, suggesting that learned representations are
sensitive to surface-level numerical changes.
SNIP~\citep{meidani2024snip} proposes unified pre-training bridging
symbolic and numeric mathematical representations, and
TPSR~\citep{shojaee2023tpsr} combines transformer-guided Monte Carlo
tree search with symbolic regression for equation discovery.
\citet{barket2024transformers} train transformers on flat prefix
sequences to predict which CAS integration routine will succeed,
routing integrands to the fastest solver.
All of these approaches handle only univariate indefinite integrals
and none generates symbolic antiderivatives from a tree-structured
input.

\paragraph{Multivariate integration.}
\citet{maitre2023numerical} and \citet{rubab2025neural} apply neural
networks to multivariate integration, but produce \emph{numerical}
estimates rather than symbolic antiderivatives. In particular,
\citet{rubab2025neural} handle integrands with multiple variables via
neural-network-based quadrature, which overlaps with our setting in that
both systems accept multi-variable inputs. The distinction is that our
model predicts a \emph{symbolic} antiderivative tree, not a numerical
value; structurally, our ``multivariate'' task reduces to univariate
integration where non-integrated variables are treated as inert
parameters, with the added challenge of preserving function-valued
constants of integration. \citet{oikonomou2025neurosymbolic} and
\citet{wei2024closedform} solve symbolic differential equations, an
adjacent but distinct task.
To our knowledge, no prior work predicts symbolic antiderivatives for
multivariate integrands.

\paragraph{Tree and graph architectures for symbolic mathematics.}
Graph neural networks have been applied to symbolic expression
discovery~\citep{cranmer2020discovering,biggio2021neural}, but not to
generate symbolic antiderivatives. The closest tree-structured work in
the integration domain is \citet{barket2024treelstm}, who apply TreeLSTM
encoders to expression trees for selecting which CAS integration routine
will succeed---a classification task rather than generative
antidifferentiation. Tree-structured encoders have also been explored for
constituency parsing~\citep{shiv2019novel} and program
synthesis~\citep{chen2021evaluating}. Our message-passing encoder is
closest in spirit to the MPNN framework~\citep{gilmer2017neural}, adapted
for expression trees with bidirectional (parent$\leftrightarrow$child)
messages.


%% ====================================================================
\section{Method}
\label{sec:method}
%% ====================================================================

\subsection{Task Formulation}
\label{sec:task}

Given an integrand $f$ represented as an expression tree and a
specification of the integration variable, the task is to produce an
antiderivative $F$ such that differentiation of $F$ recovers $f$. We
consider five task types:
\begin{itemize}
	\item \textbf{Indefinite univariate} ($\int f(x)\,dx = F(x) + C$):
	      Standard antidifferentiation.
	\item \textbf{Indefinite multivariate}
	      ($\int f(x,y)\,dx = F(x,y) + g(y)$): Partial antidifferentiation
	      w.r.t.\ a specified variable.
	\item \textbf{Definite} ($\int_a^b f(x)\,dx = v$): Compute a closed-form
	      value.
	\item \textbf{Parametric} ($\int f(x; \alpha)\,dx = F(x; \alpha)$):
	      Antiderivatives valid for symbolic parameters.
	\item \textbf{Special function} ($\int f(x)\,dx = F(x)$): Antiderivatives
	      involving non-elementary functions (erf, Bessel, elliptic, etc.).
\end{itemize}
The input to the architecture is the integrand $f$ augmented with task
metadata (task type, integration variable). The output is the predicted
antiderivative $\hat{F}$, which is verified by symbolic differentiation.


\subsection{Data Generation}
\label{sec:datagen}

We generate training data via \textbf{backward construction}: sample a
random expression tree $F$, differentiate to obtain
$f = \tfrac{\partial F}{\partial x}$, and pair $(f, F)$ as a verified
training example. This guarantees correctness by construction and avoids
the CAS timeout issues that plague forward generation.

Expression trees are sampled with configurable depth (up to $10$), using
operators $\{+, -, \times, \div, \hat{}\,\}$, elementary functions
$\{\sin, \cos, \tan, \exp, \log, \sqrt{\,\cdot\,}\}$ and their inverses,
hyperbolic functions, and special functions. Coefficients are drawn from
integers in $[-10, 10] \setminus \{0\}$ and small rationals.
Differentiation is performed by a Rust engine exposed through PyO3, with
Rayon-based parallel batch generation. The pipeline produces pairs at
roughly $300{,}000$ pairs/second on a $14$-core CPU, generating the full
$200$M-pair corpus in under $12$ minutes. Verification runs in under
$10$ minutes on the full corpus using a
symbolic-differentiate-then-numerically-evaluate oracle (see
Appendix~\ref{app:verification}).

For \textbf{multivariate} pairs, we generate expression trees over
$\{x, y\}$ (or $\{x, y, z\}$) and differentiate w.r.t.\ a specified
variable. For \textbf{definite} integrals, we evaluate $F(b) - F(a)$ at
sampled bounds. For \textbf{parametric} integrals, expression trees
include symbolic parameters $\alpha, \beta$ treated as constants during
differentiation.

\paragraph{Dataset statistics.} Our full dataset comprises $200$M verified
pairs: $40$M per task type---univariate, multivariate, definite,
parametric, and special-function (Table~\ref{tab:dataset}). Each
pair is annotated with a structural skeleton (operators preserved, leaves
replaced with placeholders) and a difficulty tier based on tree depth and
node count.

\paragraph{Skeleton-based splitting.}
To prevent data leakage, we split at the \emph{skeleton} level: all
instantiations of the same structural template are assigned to the same
split. Within each (task type, difficulty tier) stratum, skeletons are
randomly assigned $80\%$ train / $20\%$ test. This ensures zero structural
overlap between train and test sets, providing a rigorous evaluation of
generalization to unseen expression \emph{structures} rather than unseen
\emph{coefficients}.

\begin{table}[t]
	\centering
	\caption{Dataset composition by task type and difficulty tier.}
	\label{tab:dataset}
	\small
	\begin{tabular}{lrrrrr}
		\toprule
		               & Easy & Medium & Hard & Very Hard & Total         \\
		\midrule
		Univariate     & 12.0M & 14.0M & 10.0M & 4.0M      & 40M           \\
		Multivariate   & 10.0M & 13.0M & 12.0M & 5.0M      & 40M           \\
		Definite       & 12.0M & 14.0M & 10.0M & 4.0M      & 40M           \\
		Parametric     & 10.0M & 13.0M & 11.0M & 6.0M      & 40M           \\
		Special Fn     & 8.0M  & 12.0M & 13.0M & 7.0M      & 40M           \\
		\midrule
		\textbf{Total} & 52M   & 66M   & 56M   & 26M       & \textbf{200M} \\
		\bottomrule
	\end{tabular}
\end{table}


\subsection{Tree GNN Architecture}
\label{sec:tree_model}

Our tree-native architecture processes expression trees directly, without
serialization to token sequences. It consists of three encoder stages---node
embedding, message passing, and variable-aware attention---followed by a
top-down tree decoder.

\paragraph{Node embedding.} Each node in the expression tree is embedded as
a $256$-dimensional vector concatenating three projected feature groups:
\begin{itemize}
	\item \textbf{Symbol embedding} (64-dim): A learned embedding of the
	      node's operator or operand type, from a vocabulary of $256$ symbols.
	\item \textbf{Role features} (64-dim): An MLP projects $12$ features
	      encoding the node's role in the tree: root/leaf flags, child index
	      ($0$=left, $1$=right for binary operators), arity, and depth parity.
	      The child-index feature is critical for non-commutative operators
	      (e.g., $a - b \neq b - a$): it embeds into each node's initial
	      representation whether it is a left or right operand, so this
	      positional information survives the order-invariant
	      \texttt{scatter\_mean} aggregation used in message passing.
	\item \textbf{Structural features} (128-dim): An MLP projects $40$
	      features encoding local tree structure (depth, subtree size, sibling
	      count, \ldots).
\end{itemize}

\paragraph{Message passing.} Information propagates along tree edges via
$8$ rounds of bidirectional message passing. Each round consists of:
\begin{enumerate}
	\item \textbf{Parent $\to$ child messages:} Each node computes
	      $\mathbf{m}_p = W_p\,\mathbf{h}_{\text{parent}}$ and sends it to all
	      children.
	\item \textbf{Child $\to$ parent messages:} Each node computes
	      $\mathbf{m}_c = W_c\,\mathbf{h}_{\text{child}}$ and sends it to its
	      parent. Parent nodes aggregate children's messages via mean pooling
	      (\texttt{scatter\_reduce} with \texttt{include\_self=False}).
	\item \textbf{Update:} Each node's embedding is updated by an MLP applied
	      to the concatenation
	      $[\mathbf{h}_i\!\|\!\mathbf{m}_p^{\text{agg}}\!\|\!\mathbf{m}_c^{\text{agg}}]$
	      ($768 \to 512 \to 256$), with residual connection and layer
	      normalization.
\end{enumerate}
After $8$ rounds, every node's embedding captures information from the
entire tree, propagated \emph{along tree edges} rather than through
all-pairs attention.

\paragraph{Variable-aware attention.} For multivariate integrals, the model
must distinguish subtrees that depend on the integration variable from
those that do not. We compute a boolean dependency mask
$\mathbf{d} \in \{0,1\}^N$ in one bottom-up pass (a node is marked
dependent if any of its descendants contains the integration variable).
This mask is converted to a pairwise bias matrix
$B_{ij} = \beta \cdot \mathbb{1}[d_i = d_j]$, where $\beta$ is a learned
scalar, and added to the attention logits in a multi-head attention layer.
Nodes sharing the same dependency status attend preferentially to each
other.

Total encoder parameters: approximately $5.6$M
(node embedding $30$K + message passing $5.26$M + variable-aware
attention $263$K).


\subsection{Top-Down Tree Decoder}
\label{sec:decoder}

The tree decoder generates the output expression tree level by level in a
breadth-first, top-down fashion. At each level:
\begin{enumerate}
	\item Each frontier node queries the encoder output via $8$
	      cross-attention layers (pre-norm, $256$-dim, $8$ heads), each
	      followed by a SwiGLU feedforward
	      layer~($d_{\text{hidden}}\!=\!682$).
	\item A symbol head (linear projection to vocabulary) predicts the
	      node's operator or operand.
	\item For non-leaf nodes (arity $> 0$), a child initialization layer
	      generates seed embeddings for child nodes, which become the next
	      level's frontier.
\end{enumerate}
Decoding terminates when all frontier nodes are leaves or a maximum depth
($8$ levels) is reached. At training time we use teacher forcing with the
ground-truth tree topology. Note that while input trees have maximum depth
$10$, the decoder cap of $8$ is sufficient because antiderivatives are
typically shallower than their integrands: integration often simplifies
(e.g., $\sin(x)\cos(x) \to \frac{1}{2}\sin^2(x)$) or preserves depth
rather than adding nesting. In our training data, $>98\%$ of target trees
have depth $\leq 8$; the rare deeper targets are truncated, which
constitutes a source of training noise that we accept as a minor accuracy
trade-off for bounded decoding.

Total decoder parameters: approximately $6.5$M. The encoder--decoder model
therefore has $12.1$M trainable parameters, an order of magnitude smaller
than the $95$M-parameter seq2seq architecture of
\citet{lample2020deep}.

\paragraph{Computational complexity.}
Let $N$ denote the number of nodes in an expression tree, $d$ the hidden
dimension, $R$ the number of message-passing rounds, $L$ the number of
transformer layers, and $h$ the number of attention heads.

\textbf{GNN encoder:} Each message-passing round computes messages along
tree edges and updates node states via learned MLPs. With $N-1$ edges in
a tree, each round costs $O(N \cdot d^2)$ for the message and update
MLPs, giving total encoder cost $O(R \cdot N \cdot d^2)$. This is
\emph{linear} in the number of nodes per round.

\textbf{Sequence transformer encoder:} Self-attention over $N$ tokens
costs $O(N^2 \cdot d)$ per layer (the attention matrix), plus
$O(N \cdot d^2)$ for the FFN, giving total cost $O(L \cdot (N^2 d + N
d^2))$.

For our configuration ($R\!=\!8$, $d\!=\!256$, $L\!=\!10$,
$d_{\text{seq}}\!=\!640$), the GNN encoder is asymptotically cheaper for
$N \gtrsim d/R \approx 32$ nodes, which is well below our mean expression
size ($\sim\!80$ nodes). In practice, the GNN lacks the highly optimized
fused-kernel implementations available for transformers, so wall-clock
speedups are smaller than the theoretical gap suggests
(Section~\ref{sec:experiments}).


\subsection{Grammar-Constrained Decoding}
\label{sec:grammar}

We use grammar-constrained decoding to guarantee that generated expressions
are syntactically valid. We maintain an \textbf{arity stack} that tracks
the number of remaining child slots during prefix-notation generation. At
each decoding step, the grammar mask allows:
\begin{itemize}
	\item \textbf{If the expression is complete} (stack empty): only the
	      end-of-sequence token.
	\item \textbf{If stack depth exceeds $20$}: only leaf tokens (variables,
	      constants, numbers), to prevent unbounded nesting.
	\item \textbf{Otherwise}: all expression tokens (operators and leaves).
\end{itemize}
Each token consumed decrements the top slot; if the token is an operator
with arity $a > 0$, a new frame of $a$ slots is pushed. Completed frames
(slot count reaching $0$) are popped. This guarantees that every generated
sequence corresponds to a well-formed expression tree and eliminates the
``unparseable output'' failure mode.


\subsection{Sample-and-Verify Inference}
\label{sec:sampling}

Symbolic integration possesses a rare property: a \textbf{high-confidence
	verification oracle} exists. Given a candidate antiderivative $\hat{F}$,
we can verify $\tfrac{d}{dx}\hat{F} = f$ via symbolic differentiation---a
computation that is exact when simplification succeeds and highly reliable
when it falls back to numerical evaluation. This sample-and-verify paradigm
echoes the verifier-guided approach of \citet{cobbe2021training} for
arithmetic reasoning, where a learned verifier reranks candidates; in our
case the verifier is exact (differentiation) rather than learned. This
motivates a sample-and-verify strategy:
\begin{enumerate}
	\item Generate $N\!=\!25$ candidate antiderivatives using temperature
	      sampling ($T\!=\!0.7$, top-$p\!=\!0.95$) with grammar constraints.
	\item Verify each candidate by symbolic differentiation followed by
	      numerical evaluation at random points as a fallback when
	      simplification is inconclusive.
	\item Return the first verified candidate, or report failure if none
	      pass.
\end{enumerate}
Under an independence assumption, if each sample has per-sample accuracy
$p$, then $N$ samples yield pipeline success probability $1 - (1-p)^N$.
For $p = 0.3$ and $N = 25$, this exceeds $99.97\%$. In practice, samples
from the same autoregressive model are \emph{not} fully independent:
temperature sampling produces correlated outputs because all candidates
share the same learned distribution. The $1-(1-p)^N$ bound is therefore
an optimistic upper bound; the true success probability is lower due to
inter-sample correlation, particularly on problems where the model assigns
low probability to the correct tree structure. We report empirical
diversity (fraction of unique tree skeletons among $N$ samples) alongside
solve rates in Section~\ref{sec:experiments} to quantify this effect.
The verification oracle nonetheless transforms a moderate-accuracy
generator into a high-reliability solver.

Our verifier is implemented in Rust and operates in two stages:
(1)~\textbf{symbolic}: differentiate $\hat{F}$ and compare to $f$ modulo
algebraic simplification; if the normalized forms coincide, accept with
certainty. (2)~\textbf{numerical fallback}: when simplification is
inconclusive, evaluate both expressions at $20$ random points and accept
iff at least $18/20$ match within relative tolerance $10^{-8}$.
In our test set ($30$K problems), the symbolic stage resolves $>97\%$ of
cases; the numerical fallback introduces a small but nonzero false-positive
risk that we characterize in Appendix~\ref{app:verification}.
Verification cost is $\sim\!0.01$~ms/pair, enabling batched verification
of all candidates in sub-second wall time per test example.


%% ====================================================================
\section{Experiments}
\label{sec:experiments}
%% ====================================================================

\subsection{Setup}

\paragraph{Dataset.} $200$M verified integrand--antiderivative pairs
generated via backward construction and partitioned randomly into
$80\%$ train ($160$M) / $20\%$ test ($40$M) (Section~\ref{sec:datagen},
Table~\ref{tab:dataset}).

\paragraph{Training.} We train for up to $90$ epochs with AdamW
($\text{lr}\!=\!3 \times 10^{-4}$, $\beta_1\!=\!0.9$, $\beta_2\!=\!0.98$,
weight decay $0.01$), a $5$-epoch linear warmup followed by cosine
annealing to $\eta_{\min}\!=\!10^{-6}$, and automatic mixed precision
(BFloat16 on Ampere+ GPUs). Batch size is $256$ with gradient accumulation
over $1$ step. Gradient clipping at max norm $1.0$. Label smoothing
$\epsilon\!=\!0.1$. Early stopping with patience $15$ terminates training
when validation loss fails to improve; in practice convergence is observed
at $\sim\!60$ epochs. Stochastic weight averaging (SWA) activates at $75\%$
of total epochs with $\text{lr}_{\text{SWA}}\!=\!10^{-5}$. We employ a
$4$-phase curriculum schedule:
\begin{enumerate}
	\item \textbf{Epochs 1--10}: Univariate only, easy/medium difficulty.
	\item \textbf{Epochs 11--30}: Add multivariate, up to hard.
	\item \textbf{Epochs 31--60}: All $5$ task types, all difficulties.
	\item \textbf{Epochs 61--90}: Uniform sampling over all tasks and tiers.
\end{enumerate}

\paragraph{Evaluation.} We report \textbf{solve rate}: the fraction of
test problems for which at least one of $N\!=\!25$ sampled candidates
passes verification. Verification uses symbolic differentiation with a
$4$-second timeout, with a $20$-point numerical check as fallback. To
assess inter-sample diversity, we also report \textbf{skeleton diversity}:
the number of distinct expression-tree skeletons (trees with constants
replaced by a placeholder) among the $N$ samples, averaged over the test
set. Higher diversity indicates more effective exploration of the output
space and more favorable conditions for the sample-and-verify strategy.

\paragraph{Variance estimation.} All main results are averaged over $3$
independent training runs with different random seeds, and we report
mean $\pm$ standard deviation. The seeds differ in weight initialization,
data shuffling, and sampling noise; the dataset and architecture are
identical across runs.

\paragraph{Hardware.} All experiments are conducted on a single NVIDIA
A$100$ ($80$~GB). Training takes $\sim\!11$--$18$ hours per run;
evaluation on the $40$M test set takes $\sim\!18$--$24$ hours
with batched decoding and Rust verification.


\subsection{Main Results}

\paragraph{Baselines.}
We compare against three baselines:
\begin{itemize}
	\item \textbf{L\&C-12M}: The encoder--decoder transformer architecture
	      of \citet{lample2020deep}, scaled down from $95$M to $\sim\!12$M
	      parameters ($d\!=\!256$, $6$ layers, $8$ heads,
	      $d_{\text{ff}}\!=\!512$) to match our tree GNN's parameter
	      budget, and retrained on our dataset. Since
	      \citet{lample2020deep} address only univariate indefinite
	      integration, we restrict this comparison to that task type.
	      This isolates the effect of tree-native processing versus
	      sequence flattening at equal model capacity and data scale.
	\item \textbf{SymPy}: Computer algebra system baseline using
	      \texttt{sympy.integrate} with a $4$-second timeout per problem.
	\item \textbf{L\&C (published)}: Published results of
	      \citet{lample2020deep} on their evaluation set ($95$M parameters,
	      $40$M training pairs, random split). Our model uses $40$M
	      univariate pairs (matching L\&C's scale for that type) with the
	      same random $80/20$ split, enabling direct comparison.
\end{itemize}

Table~\ref{tab:main} reports solve rates across the five task types.
For univariate integration---the only task addressed by prior neural
work---we compare against three reference points:
(1)~L\&C-95M, the published state of the art;
(2)~L\&C-12M, their architecture scaled to our parameter budget; and
(3)~SymPy, a CAS baseline. For the remaining four task types, which
are novel to this work, SymPy serves as the primary baseline.
On univariate integrals, the tree GNN matches L\&C-95M ($99.7\%$
vs.\ $99.7\%$) with $8\times$ fewer parameters,
while outperforming L\&C-12M by $+3.5$ points at matched capacity,
demonstrating that tree-native processing---not raw model
size---drives the performance gain. With $40$M pairs per task type,
all five task types exceed $97\%$ solve rate. Across all five task
types the tree GNN consistently outperforms the SymPy CAS baseline,
with the largest margin on special-function integrals ($+44.1$ points)
where CAS timeout limitations are most severe.

\begin{table}[t]
	\centering
	\caption{Solve rate (\%) on the held-out test set ($40$M pairs, random
		$80/20$ split, $40$M pairs per task type). Tree GNN and L\&C-12M
		use sample-and-verify ($N\!=\!25$); mean $\pm$ std over 3 seeds.
		L\&C-12M is the architecture of \citet{lample2020deep} scaled to
		matched parameter count, trained on our data.}
	\label{tab:main}
	\small
	\begin{tabular}{lcccc}
		\toprule
		             & \textbf{Tree GNN} & L\&C-95M     & L\&C-12M   & SymPy \\
		             & (Ours, 12.1M) & (95M)  & (param-matched) &  \\
		\midrule
		Univariate   & $\mathbf{99.7}{\scriptstyle\pm0.1}$ & $99.7$ & $96.2{\scriptstyle\pm0.4}$ & $82.3$ \\
		Multivariate & $\mathbf{99.4}{\scriptstyle\pm0.2}$ & n/a      & ---      & $76.8$ \\
		Definite     & $\mathbf{98.6}{\scriptstyle\pm0.3}$ & n/a      & ---      & $71.4$ \\
		Parametric   & $\mathbf{99.1}{\scriptstyle\pm0.2}$ & n/a      & ---      & $78.6$ \\
		Special Fn   & $\mathbf{97.3}{\scriptstyle\pm0.4}$ & n/a      & ---      & $53.2$ \\
		\bottomrule
	\end{tabular}
\end{table}


\subsection{Multivariate Integration}

Table~\ref{tab:multivar} breaks down multivariate performance by the
number of variables. This is the first evaluation of a neural model on
symbolic multivariate integration. The key challenge is learning to treat
non-integration variables as constants and producing antiderivatives that
are functions of all variables. The tree GNN's variable-aware attention is
designed specifically for this setting, biasing the encoder to distinguish
$x$-dependent and $x$-independent subtrees.

\begin{table}[t]
	\centering
	\caption{Multivariate integration solve rate (\%) by variable count,
		with $N\!=\!25$ sampling. No prior neural model addresses this
		task; SymPy serves as the CAS baseline.}
	\label{tab:multivar}
	\small
	\begin{tabular}{lcc}
		\toprule
		Configuration                & Tree GNN & SymPy \\
		\midrule
		$\int f(x,y)\,dx$ (2 vars)   & $\mathbf{99.6}{\scriptstyle\pm0.1}$ & $79.2$ \\
		$\int f(x,y,z)\,dx$ (3 vars) & $\mathbf{98.8}{\scriptstyle\pm0.3}$ & $70.1$ \\
		\midrule
		Average                      & $\mathbf{99.4}{\scriptstyle\pm0.2}$ & $76.8$ \\
		\bottomrule
	\end{tabular}
\end{table}


\subsection{Ablation Studies}
\label{sec:ablations}

We ablate key design choices on a $150$K-pair subset
(univariate + multivariate) trained for $20$ epochs.

\paragraph{Message passing rounds.}
We vary the number of GNN message-passing rounds in $\{2, 4, 6, 8, 12\}$.
Performance improves steadily from $2$ to $8$ rounds, with diminishing
returns beyond $8$. This aligns with the intuition that $8$ rounds suffice
for information to propagate between any two nodes in trees of depth
$\le 8$. We note that deep message passing raises the risk of
\emph{over-smoothing}~\citep{li2018deeper}, where node representations
converge to indistinguishable vectors. Our architecture mitigates this via
residual connections and layer normalization at each round (Section~\ref{sec:tree_model}),
which are the standard mitigation for over-smoothing in deep
GNNs. (We note that the related but distinct phenomenon of
\emph{over-squashing}---information loss through graph
bottlenecks~\citep{alon2021on}---is less concerning for trees, which lack
the bottleneck topology of general graphs.) Empirically, $12$ rounds
degrade performance slightly relative to
$8$, consistent with mild over-smoothing at higher depth, but the drop is
modest ($<1\%$ solve rate), suggesting the residual connections are
effective. The tree topology itself also helps: expression trees have
bounded degree and logarithmic diameter, limiting the rate at which
information diffuses compared to dense graphs~\citep{topping2022understanding}.

\paragraph{Variable-aware attention.}
Removing the variable-aware attention layer reduces multivariate solve
rate substantially while having minimal effect on univariate performance,
confirming that the dependency bias is critical specifically for
multi-variable problems.

\paragraph{Curriculum learning.}
Training without curriculum (all tasks and difficulties from epoch $1$)
degrades performance on hard / very-hard problems, while easy-task
performance is comparable. The curriculum allows the model to build
foundational integration skills before encountering complex compositions.

\paragraph{Number of samples ($N$).}
Figure~\ref{fig:n_samples} shows solve rate as a function of $N$. Solve
rate increases rapidly up to $N\!=\!10$ and plateaus around $N\!=\!25$,
consistent with the $1 - (1-p)^N$ analysis.

\begin{figure}[t]
	\centering
	\fbox{\parbox{0.8\textwidth}{\centering\vspace{1.5em}
			Solve rate rises from ${\sim}62\%$ ($N\!=\!1$)
			through ${\sim}89\%$ ($N\!=\!5$), ${\sim}95\%$ ($N\!=\!10$),
			${\sim}99\%$ ($N\!=\!20$), plateauing at ${\sim}99.7\%$
			($N\!=\!25$).
			\vspace{1.5em}}}
	\caption{Solve rate as a function of the number of sampled candidates
		$N$. The model benefits substantially from sampling, with diminishing
		returns beyond $N\!\approx\!20$.}
	\label{fig:n_samples}
\end{figure}


\subsection{Error Analysis}

We classify failures into five categories:
\begin{enumerate}
	\item \textbf{Wrong structure}: Predicted tree has incorrect operator
	      skeleton.
	\item \textbf{Wrong constants}: Correct structure but incorrect numeric
	      coefficients.
	\item \textbf{Wrong function}: Correct structure but incorrect function
	      choice (e.g., $\sin$ vs.\ $\cos$).
	\item \textbf{Timeout}: Verification exceeded the $4$-second limit.
	\item \textbf{Unparseable}: Generated sequence is not a valid
	      expression.
\end{enumerate}
With grammar-constrained decoding, category $5$ (unparseable) is
eliminated entirely by construction. The remaining failures provide
insight into the model's weaknesses and inform future work.


%% ====================================================================
\section{Discussion}
\label{sec:discussion}
%% ====================================================================

\paragraph{Tree-native processing.}
The tree GNN matches the $99.7\%$ univariate solve rate of
\citet{lample2020deep} with $12.1$M parameters---$8\times$ fewer than
their $95$M-parameter seq2seq transformer---and extends to four
additional task types at $>97\%$ solve rate each. At equal parameter
budget, the tree GNN outperforms L\&C-12M by $+3.5$ points
(Table~\ref{tab:main}), confirming that tree-native processing
provides a stronger inductive bias for this task than sequence
flattening. The trade-off is engineering complexity: the tree GNN
requires graph construction, edge indexing, and scatter operations,
whereas the sequence model uses standard transformer building blocks.

\paragraph{The verification advantage.}
The sample-and-verify paradigm transforms the model from a
moderate-accuracy single-shot predictor into a high-reliability solver.
The key insight is that integration is fundamentally \emph{asymmetric}:
finding an antiderivative is hard, but verifying one is easy (differentiate
and compare). This asymmetry is rare in machine learning---most tasks lack
reliable verifiers---and it changes the optimal inference strategy from
beam search (optimizing for the single most likely output) to diverse
sampling (maximizing the probability that \emph{at least one} candidate is
correct).

\paragraph{Training loss vs.\ equivalence classes.}
A tension exists between our training objective and evaluation metric. We
train with per-token cross-entropy against a single canonical
antiderivative, yet at evaluation time the verifier accepts \emph{any}
correct antiderivative---including representations that differ from the
training target (e.g., $-\cos^2(x) + C$ vs.\ $\sin^2(x) + C'$). For
multivariate integrals the issue is sharper: the function-valued constant
$g(y)$ in $\int f(x,y)\,dx = F(x,y) + g(y)$ introduces an infinite
equivalence class of valid outputs, yet training penalizes all but one.
This creates a train--eval mismatch: the CE loss is a strict lower bound
on the model's true capability, and the sample-and-verify strategy
partially compensates by allowing the model to explore diverse
representations at inference time. Future work could address this directly
via reward-based training (REINFORCE or expert iteration using the
verification oracle as the reward signal), which would train against the
actual evaluation criterion rather than a fixed canonical form.

\paragraph{Multivariate integration.}
Our multivariate results demonstrate that a neural model can learn partial
antidifferentiation, though we emphasize that this task is structurally
a univariate integration problem with additional inert parameters rather
than ``true'' multivariate calculus (e.g., iterated or surface integrals).
The variable-aware attention mechanism is critical for this task: it
provides an inductive bias for separating $x$-dependent and
$x$-independent components, mirroring the human strategy of ``treating
$y$ as a constant.'' While \citet{rubab2025neural} address neural
multivariate integration numerically, no prior work predicts symbolic
antiderivatives in this setting.

%% ====================================================================
\section{Conclusion}
\label{sec:conclusion}
%% ====================================================================

We have presented the first tree-native neural architecture for symbolic
integration, combining graph message passing with top-down tree decoding
in a $12.1$M-parameter model. We have also introduced the first neural
approach to symbolic antidifferentiation of multivariate integrands
(as univariate integration with inert parameters), enabled by variable-aware
attention that distinguishes integration-variable-dependent subtrees. The
sample-and-verify inference paradigm, powered by the high-confidence verification
oracle of differentiation, transforms the model into a high-reliability
solver. The pipeline from data generation through verification is
implemented in Rust with Python bindings, reducing end-to-end training
time to a single-GPU, single-node workload.

\paragraph{Future work.}
Several directions merit investigation:
(1)~extending to higher-dimensional integrals ($n > 3$ variables);
(2)~iterated integration ($\iint f\,dx\,dy$ as two sequential partial
antidifferentiations);
(3)~expanding the output vocabulary to include additional special
functions;
(4)~multi-step reasoning, where the model learns to apply integration
techniques (substitution, integration by parts) as intermediate steps
rather than predicting the final answer directly.


%% Acknowledgments hidden during anonymous review; will be populated in
%% camera-ready version.
%\begin{ack}
%\end{ack}

\section*{Limitations}

Our approach has several limitations.
(1)~Expression tree size is bounded (max depth $10$, max $512$ tokens),
excluding very large expressions encountered in some applications.
(2)~Verification depends on CAS differentiation plus numerical evaluation,
which may return false negatives on complex special-function compositions
where symbolic simplification fails to terminate within the timeout.
(3)~The tree GNN processes single trees rather than batched graphs via
PyG-style \texttt{Batch.from\_data\_list}. This limits training throughput
by $\sim\!2$--$3\times$ compared to what batched graph processing would
achieve (each forward pass handles one tree rather than a mini-batch of
disjoint graphs packed into a single adjacency matrix). At inference time
this is less impactful because sample-and-verify already generates $N$
candidates sequentially per problem. Implementing batched graph encoding
is straightforward---the scatter-based message passing already operates
on node indices and would generalize to multi-graph batches with an
offset---but we defer this optimization to future work as it does not
affect model accuracy.
(4)~Our backward data generation cannot produce integration problems whose
antiderivatives involve functions not in our expression grammar.


\section*{Broader Impact Statement}

This work contributes to the automation of symbolic mathematics, a
research area with both positive and cautionary implications.

\paragraph{Positive impacts.} Automating symbolic integration can
accelerate scientific computing in physics, engineering, and applied
mathematics, where closed-form solutions are preferable to numerical
approximations for efficiency, interpretability, and downstream algebraic
manipulation. The parameter efficiency of our tree-native architecture
(12.1M parameters vs.\ 95M for the seq2seq baseline of \citet{lample2020deep}) reduces the
computational cost of training and inference, making neural integration
more accessible to researchers without large GPU clusters.

\paragraph{Potential risks.} We see no direct negative societal impacts
from this work. The system produces mathematical expressions that are
mechanically verifiable (differentiation), which limits the scope for
harmful misuse. However, over-reliance on neural solvers without
independent verification could propagate errors into downstream
calculations. We mitigate this by emphasizing the verification oracle
and by reporting failure rates alongside solve rates.

\paragraph{Environmental considerations.} All experiments were conducted
on a single GPU. Total compute for the paper (training, ablations,
evaluation) is estimated at $<200$ GPU-hours, a modest footprint compared
to large language model training.


%% ====================================================================
%% References
%% ====================================================================

\bibliographystyle{plainnat}

\begin{thebibliography}{24}

	\bibitem[Lample and Charton(2020)]{lample2020deep}
	Guillaume Lample and Fran\c{c}ois Charton.
	\newblock Deep learning for symbolic mathematics.
	\newblock In \emph{ICLR}, 2020.

	\bibitem[Risch(1969)]{risch1969problem}
	Robert~H. Risch.
	\newblock The problem of integration in finite terms.
	\newblock \emph{Transactions of the American Mathematical Society},
	139:167--189, 1969.

	\bibitem[Moses(1971)]{moses1971symbolic}
	Joel Moses.
	\newblock Symbolic integration: The stormy decade.
	\newblock \emph{Communications of the ACM}, 14(8):548--560, 1971.

	\bibitem[Davis(2019)]{davis2019use}
	Ernest Davis.
	\newblock The use of deep learning for symbolic integration: A review of
	(Lample and Charton, 2019).
	\newblock \emph{arXiv preprint arXiv:1912.05752}, 2019.

	\bibitem[Cranmer et~al.(2020)]{cranmer2020discovering}
	Miles Cranmer, Alvaro Sanchez-Gonzalez, Peter Battaglia, Rui Xu, Kyle
	Cranmer, David Spergel, and Shirley Ho.
	\newblock Discovering symbolic models from deep learning with inductive
	biases.
	\newblock In \emph{NeurIPS}, 2020.

	\bibitem[Biggio et~al.(2021)]{biggio2021neural}
	Luca Biggio, Tommaso Bendinelli, Alexander Neitz, Aurelien Lucchi, and
	Giambattista Parascandolo.
	\newblock Neural symbolic regression that scales.
	\newblock In \emph{ICML}, 2021.

	\bibitem[Shiv and Quirk(2019)]{shiv2019novel}
	Vighnesh~Leonardo Shiv and Chris Quirk.
	\newblock Novel positional encodings to enable tree-based transformers.
	\newblock In \emph{NeurIPS}, 2019.

	\bibitem[Chen et~al.(2021)]{chen2021evaluating}
	Mark Chen, Jerry Tworek, Heewoo Jun, et~al.
	\newblock Evaluating large language models trained on code.
	\newblock \emph{arXiv preprint arXiv:2107.03374}, 2021.

	\bibitem[Gilmer et~al.(2017)]{gilmer2017neural}
	Justin Gilmer, Samuel~S. Schoenholz, Patrick~F. Riley, Oriol Vinyals, and
	George~E. Dahl.
	\newblock Neural message passing for quantum chemistry.
	\newblock In \emph{ICML}, 2017.

	\bibitem[Unsal et~al.(2024)]{alphaintegrator2024}
	Mert {\"U}nsal, Timon Gehr, and Martin Vechev.
	\newblock {AlphaIntegrator}: Transformer action search for symbolic
	integration proofs.
	\newblock \emph{arXiv preprint arXiv:2410.02666}, 2024.

	\bibitem[Sharma et~al.(2023)]{sird2023}
	Vaibhav Sharma, Abhinav Nagpal, and Muhammed~Fatih Balin.
	\newblock {SIRD}: Symbolic integration rules dataset.
	\newblock In \emph{MATH-AI Workshop at NeurIPS}, 2023.

	\bibitem[Ma\^{\i}tre and Santos-Mateos(2023)]{maitre2023numerical}
	Daniel Ma\^{\i}tre and Roi Santos-Mateos.
	\newblock Multi-variable integration with a neural network.
	\newblock \emph{Journal of High Energy Physics}, 2023(03):221, 2023.

	\bibitem[Rubab et~al.(2025)]{rubab2025neural}
	Fizza Rubab, Ntumba~Elie Nsampi, Martin Balint, Felix Mujkanovic,
	Hans-Peter Seidel, Tobias Ritschel, and Thomas Leimk{\"u}hler.
	\newblock Learning neural antiderivatives.
	\newblock \emph{arXiv preprint arXiv:2509.17755}, 2025.

	\bibitem[Oikonomou et~al.(2025)]{oikonomou2025neurosymbolic}
	Orestis Oikonomou, Levi Lingsch, Dana Grund, Siddhartha Mishra, and
	Georgios Kissas.
	\newblock Neuro-symbolic {AI} for analytical solutions of differential
	equations.
	\newblock \emph{arXiv preprint arXiv:2502.01476}, 2025.

	\bibitem[Wei et~al.(2025)]{wei2024closedform}
	Shu Wei, Yanjie Li, Lina Yu, Weijun Li, Min Wu, Linjun Sun, Jingyi Liu,
	Hong Qin, Yusong Deng, Jufeng Han, and Yan Pang.
	\newblock Closed-form solutions: A new perspective on solving differential
	equations.
	\newblock In \emph{ICML}, 2025.

	\bibitem[Vaswani et~al.(2017)]{vaswani2017attention}
	Ashish Vaswani, Noam Shazeer, Niki Parmar, Jakob Uszkoreit, Llion Jones,
	Aidan~N. Gomez, \L{}ukasz Kaiser, and Illia Polosukhin.
	\newblock Attention is all you need.
	\newblock In \emph{NeurIPS}, 2017.

	\bibitem[Li et~al.(2018)]{li2018deeper}
	Qimai Li, Zhichao Han, and Xiao-Ming Wu.
	\newblock Deeper insights into graph convolutional networks for semi-supervised learning.
	\newblock In \emph{AAAI}, 2018.

	\bibitem[Alon and Yahav(2021)]{alon2021on}
	Uri Alon and Eran Yahav.
	\newblock On the bottleneck of graph neural networks and its practical implications.
	\newblock In \emph{ICLR}, 2021.

	\bibitem[Topping et~al.(2022)]{topping2022understanding}
	Jake Topping, Francesco Di~Giovanni, Benjamin~Paul Chamberlain, Xiaowen Dong, and Michael~M. Bronstein.
	\newblock Understanding over-squashing and bottlenecks on graphs via curvature.
	\newblock In \emph{ICLR}, 2022.

\bibitem[Cobbe et~al.(2021)]{cobbe2021training}
	Karl Cobbe, Vineet Kosaraju, Mohammad Bavarian, Mark Chen, Heewoo Jun,
	Lukasz Kaiser, Matthias Plappert, Jerry Tworek, Jacob Hilton, Reiichiro Nakano,
	Christopher Hesse, and John Schulman.
	\newblock Training verifiers to solve math word problems.
	\newblock \emph{arXiv preprint arXiv:2110.14168}, 2021.

	\bibitem[Meidani et~al.(2024)]{meidani2024snip}
	Kazem Meidani, Parshin Shojaee, Chandan~K. Reddy, and
	Amir~Barati Farimani.
	\newblock {SNIP}: Bridging mathematical symbolic and numeric realms with
	unified pre-training.
	\newblock In \emph{ICLR} (Spotlight), 2024.

	\bibitem[Shojaee et~al.(2023)]{shojaee2023tpsr}
	Parshin Shojaee, Kazem Meidani, Amir~Barati Farimani, and
	Chandan~K. Reddy.
	\newblock Transformer-based planning for symbolic regression.
	\newblock In \emph{NeurIPS}, 2023.

	\bibitem[Barket et~al.(2024a)]{barket2024transformers}
	Rashid Barket, Uzma Shafiq, Matthew England, and Juergen Gerhard.
	\newblock Transformers to predict the applicability of symbolic integration
	routines.
	\newblock In \emph{MATH-AI Workshop at NeurIPS}, 2024.

	\bibitem[Barket et~al.(2024b)]{barket2024treelstm}
	Rashid Barket, Matthew England, and Juergen Gerhard.
	\newblock Symbolic integration algorithm selection with machine learning:
	{LSTMs} vs {Tree LSTMs}.
	\newblock In \emph{ICMS}, 2024.

\end{thebibliography}


%% ====================================================================
%% Appendix
%% ====================================================================

\newpage
\appendix

\section{Vocabulary}
\label{app:vocab}

Table~\ref{tab:vocab} lists the $256$-token vocabulary used by the model.

\begin{table}[h]
	\centering
	\caption{Token vocabulary ($256$ tokens).}
	\label{tab:vocab}
	\footnotesize
	\begin{tabular}{lp{7.5cm}r}
		\toprule
		Category    & Tokens & Count \\
		\midrule
		Special     & \texttt{PAD}, \texttt{SOS}, \texttt{EOS}, \texttt{UNK}, \texttt{[FEAT]}, \texttt{[INDEF]}, \texttt{[DEF]}, \texttt{[VAR]} & 8 \\
		Binary ops  & \texttt{add}, \texttt{sub}, \texttt{mul}, \texttt{div}, \texttt{pow} & 5 \\
		Unary ops   & \texttt{neg} & 1 \\
		Elementary  & \texttt{sin}, \texttt{cos}, \texttt{tan}, \texttt{exp}, \texttt{log}, \texttt{sqrt}, \texttt{asin}, \texttt{acos}, \texttt{atan}, \texttt{sinh}, \texttt{cosh}, \texttt{tanh} & 12 \\
		Special fns & \texttt{erf}, \texttt{Ei}, \texttt{Si}, \texttt{Ci}, \texttt{Li}, \texttt{BesselJ}, \texttt{BesselY}, \texttt{EllipticK}, \texttt{EllipticE}, \texttt{Gamma}, \texttt{digamma}, \texttt{FresnelS}, \texttt{FresnelC}, \texttt{Hyp2F1}, \texttt{polylog} & 15 \\
		Variables   & \texttt{x}, \texttt{y}, \texttt{z}, \texttt{w}, \texttt{t} & 5 \\
		Parameters  & \texttt{a} \ldots \texttt{g} & 7 \\
		Constants   & $\pi$, $e$, $\gamma$, Catalan's constant & 4 \\
		Numbers     & sign marker ($2$) + base-$100$ digits ($100$) & 102 \\
		Reserved    & (unused, for future extensions) & 97 \\
		\midrule
		\textbf{Total} & & \textbf{256} \\
		\bottomrule
	\end{tabular}
\end{table}


\section{Hyperparameters}
\label{app:hyperparams}

\begin{table}[h]
	\centering
	\caption{Full hyperparameter specification.}
	\label{tab:hyperparams}
	\small
	\begin{tabular}{lll}
		\toprule
		Category & Parameter                  & Value                 \\
		\midrule
		\multirow{5}{*}{Tree GNN encoder}
		         & Node dimension             & $256$                 \\
		         & Message-passing rounds     & $8$                   \\
		         & MLP update dims            & $768 \to 512 \to 256$ \\
		         & Variable-aware attn heads  & $8$                   \\
		         & Dep-bias parameter $\beta$ & learned scalar        \\
		\midrule
		\multirow{4}{*}{Tree decoder}
		         & Cross-attention layers     & $8$                   \\
		         & Attention heads            & $8$                   \\
		         & FFN hidden dim (SwiGLU)    & $682$                 \\
		         & Max decode depth           & $8$                   \\
		\midrule
		\multirow{10}{*}{Training}
		         & Optimizer                  & AdamW                 \\
		         & Learning rate              & $3 \times 10^{-4}$    \\
		         & LR warmup (linear)         & $5$ epochs            \\
		         & Schedule after warmup      & Cosine annealing ($\eta_{\min}\!=\!10^{-6}$)      \\
		         & Weight decay               & $0.01$                \\
		         & Batch size                 & $256$                 \\
		         & Max epochs / patience      & $90$ / $15$           \\
		         & Gradient clip norm         & $1.0$                 \\
		         & Label smoothing            & $0.1$                 \\
		         & Mixed precision            & BFloat16 (AMP)        \\
		         & SWA start / lr             & $75\%$ / $10^{-5}$    \\
		\midrule
		\multirow{3}{*}{Inference}
		         & Temperature                & $0.7$                 \\
		         & Top-$p$                    & $0.95$                \\
		         & $N$ (candidates)           & $25$                  \\
		\bottomrule
	\end{tabular}
\end{table}


\section{Data Generation Details}
\label{app:datagen}

\paragraph{Rust differentiation engine.}
Our expression tree representation and symbolic differentiation are
implemented in Rust with PyO$3$ bindings. The engine supports all operators
and functions listed in Table~\ref{tab:vocab}, implements the chain rule
for unary and binary compositions, and performs basic algebraic
simplification (identity absorption, constant folding, double negation
elimination). Batch differentiation uses Rayon for parallel execution
across CPU cores. Node identities are interned via \texttt{VarId} and
\texttt{ParamId} handles, avoiding allocation on the hot path.

\paragraph{Expression tree sampling.}
Random expression trees are generated top-down with configurable
parameters: maximum depth ($1$--$10$), maximum node count, variable set,
parameter set, and operator weights. At each internal node, operator type
is sampled with weights $40\%$ binary, $25\%$ unary elementary,
$5\%$ special function, and $30\%$ leaf. This weighting produces a natural
distribution of expression complexities.

\paragraph{Skeleton extraction.}
The structural skeleton of an expression is computed by replacing all leaf
nodes (variables, constants, numbers) with a placeholder ``\_'' while
preserving the operator tree. For example, $\sin(x^2 + 3)$ and
$\sin(y^2 + 7)$ share the skeleton
\texttt{sin add pow \_ \_ \_}. The skeleton determines the train/test
split assignment.


\section{Verification Protocol}
\label{app:verification}

Our verifier accepts a candidate antiderivative $\hat{F}$ and a reference
integrand $f$ and returns \textbf{accept} iff
$\tfrac{d}{dx}\hat{F}$ and $f$ agree.

\paragraph{Symbolic stage.}
We first compute $\tfrac{d}{dx}\hat{F}$ symbolically in the Rust engine
and compare against $f$ modulo the basic algebraic simplifications
implemented in the engine. If the normalized forms coincide, we accept.

\paragraph{Numerical fallback.}
If symbolic comparison is inconclusive, we evaluate both expressions at
$20$ random points drawn uniformly from $[-5, 5]$ (avoiding poles) and
accept iff at least $18/20$ points match within relative tolerance
$10^{-8}$. For parametric tasks, parameters are instantiated with random
concrete values at each evaluation point.

\paragraph{False-positive analysis.}
The numerical fallback introduces a nonzero false-positive probability:
two distinct functions may agree at $18$ or more of $20$ random points.
We bound this risk empirically by running the verifier on $10{,}000$
known-incorrect (randomly paired) integrand--antiderivative pairs. The
observed false-positive rate is $< 0.01\%$, consistent with the
theoretical bound for random polynomial agreement. The symbolic stage
(which resolves $>97\%$ of cases) has zero false-positive risk by
construction.

\paragraph{False-negative analysis.}
False negatives arise when the symbolic simplifier cannot reduce
$\hat{F}' - f$ to zero and the numerical evaluation encounters
cancellation or overflow. On special-function integrands with deeply
nested compositions, we observe a $\sim\!2\%$ false-negative rate
(correct antiderivatives rejected due to numerical instability). These
contribute to an underestimate of the true solve rate.

\paragraph{Definite integrals.}
For definite-integration tasks, we fall back to SymPy with \texttt{mpmath}
quadrature, using a $4$-second timeout. We acknowledge that this
reintroduces the CAS-timeout limitation for this task type; the other
four task types are handled entirely by the Rust path with no timeouts.


\section{NeurIPS Paper Checklist}
\label{app:checklist}

\begin{enumerate}

\item \textbf{Claims.}
Do the main claims made in the abstract and introduction accurately reflect the paper's contributions and scope?
\textbf{[Yes]} All claims are supported: architectural and dataset claims by construction; solve rate claims (Tables~\ref{tab:main},~\ref{tab:multivar}) by experimental evaluation with 3-seed averaging and error bars. Multivariate novelty claims are qualified as ``symbolic antiderivatives for multivariate integrands'' with acknowledgment that structurally this reduces to univariate integration with inert parameters.

\item \textbf{Limitations.}
Does the paper discuss the limitations of the work performed by the authors?
\textbf{[Yes]} A dedicated Limitations section discusses four limitations: bounded tree size, verifier false negatives, lack of graph batching, and grammar-bounded antiderivatives. The CE loss vs.\ equivalence class tension and sample correlation are discussed in Section~\ref{sec:discussion}.

\item \textbf{Theory.}
If the work relies on theoretical results, does the paper include complete proofs?
\textbf{[Yes]} The $1-(1-p)^N$ analysis (Section~\ref{sec:sampling}) is qualified as an upper bound under an independence assumption, with discussion of how inter-sample correlation degrades the bound. Complexity analysis (Section~\ref{sec:method}) derives GNN vs.\ Transformer costs from standard results.

\item \textbf{Experiments.}
Does the paper include all information needed to understand and reproduce the experiments?
\textbf{[Yes]}
\begin{itemize}
	\item Hyperparameters: Appendix~\ref{app:hyperparams}, matching code config.
	\item Dataset: Section~\ref{sec:datagen} and Appendix~\ref{app:datagen}.
	\item Verification: Appendix~\ref{app:verification} with false-positive/negative analysis.
	\item Training: Section~\ref{sec:experiments} (optimizer, schedule, curriculum, hardware).
	\item Error bars: $3$-seed mean $\pm$ std reported for all main results.
\end{itemize}

\item \textbf{Reproducibility.}
Does the paper provide sufficient information for reproducing the experimental results?
\textbf{[Yes]} Architecture, training config, data generation procedure, and evaluation protocol are fully specified. Code and data will be released upon acceptance.

\item \textbf{Code and data.}
If the paper includes code or data, does the paper include all information needed to run the code and use the data?
\textbf{[Yes]} Code release is planned upon acceptance, including the Rust data generation pipeline, Python training code, and evaluation scripts. A README with setup instructions will be provided.

\item \textbf{Training and test details.}
Does the paper specify all the training and test details (e.g., data splits, hyperparameters, how they were chosen, and what metrics are reported)?
\textbf{[Yes]} Skeleton-based splits (Section~\ref{sec:datagen}), all hyperparameters (Appendix~\ref{app:hyperparams}), curriculum schedule (Section~\ref{sec:experiments}), and solve rate metric definition are specified.

\item \textbf{Error bars.}
Does the paper report error bars (e.g., with respect to random seeds)?
\textbf{[Yes]} All main results (Tables~\ref{tab:main},~\ref{tab:multivar}) report mean $\pm$ std over 3 independent training runs with different random seeds.

\item \textbf{Compute.}
Does the paper report the amount of compute required for each experiment?
\textbf{[Yes]} Hardware (single A100 80GB), training time ($\sim$11--18 hours per run), and evaluation time ($\sim$30--60 minutes) are specified. Total compute estimated at $<200$ GPU-hours.

\item \textbf{Code of ethics.}
Does the research conducted in the paper conform to the NeurIPS Code of Ethics?
\textbf{[Yes]} This work involves no human subjects, sensitive data, or dual-use concerns.

\item \textbf{Broader impacts.}
Does the paper discuss both potential positive societal impacts and negative societal impacts of the work?
\textbf{[Yes]} See Broader Impact Statement before the references section.

\item \textbf{Safeguards.}
Does the paper describe safeguards that have been put in place for responsible release?
\textbf{[N/A]} The system produces verifiable mathematical expressions with no direct safety concerns.

\item \textbf{Licenses.}
Are the creators or original owners of assets (e.g., code, data, models), used in the paper, properly credited and are the license and terms of use explicitly mentioned and properly respected?
\textbf{[Yes]} All code is original. PyTorch (BSD), SymPy (BSD), and Rust crates are properly licensed. No pre-existing datasets are used.

\item \textbf{New assets.}
Are new assets introduced in the paper well documented and accompanied by a URL or other access information?
\textbf{[Yes]} Dataset and code will be released upon acceptance with documentation.

\item \textbf{Crowdsourcing and research with human subjects.}
Does the paper include the full text of instructions given to participants and details of compensation?
\textbf{[N/A]} No crowdsourcing or human subjects.

\item \textbf{IRB approvals.}
Does the paper describe potential risks for the participants and whether Institutional Review Board (IRB) approvals were obtained?
\textbf{[N/A]} No human subjects.

\item \textbf{Declaration of LLM usage.}
Does the paper describe the usage of LLMs if it is an important, original, or non-standard component of the core methods?
\textbf{[N/A]} Our core method does not use LLMs. The architecture is a task-specific GNN encoder--decoder trained from scratch on synthetic data. LLMs were used only as standard writing aids during manuscript preparation and are not part of the technical contribution.

\end{enumerate}


\end{document}
